\emph{Översikt}
Den här modulen bygger ett enda program, från början till slut: ett
gissningsspel som spelas i flera omgångar och som till sist berättar hur
många försök du behövde i genomsnitt.  Programmet är litet, men det
kräver allt vi gått igenom samtidigt --- en behållare som samlar
resultaten medan spelet pågår, upprepningar som driver omgångarna och
gissningarna, funktioner som delar upp arbetet, och felhantering för att
användaren skriver vad som helst.  Vi utvecklar det med stegvis
förfining, och tar upp den fråga som gör ett sådant här program svårt
att visa i ett läromedel: hur man får ett program som slumpar att bete
sig likadant två gånger.

\emph{Lärandemål}
Efter denna modul ska studenten kunna:
\begin{restatable}{lo}{GuessLOAccumulate}\label{GuessLOAccumulate}%
  Använda en lista för att samla resultat under körningen, och
  sammanfatta den efteråt med \mintinline{python}{sum} och
  \mintinline{python}{len}.
\end{restatable}% Lo09, Lo04
\begin{restatable}{lo}{GuessLOValidate}\label{GuessLOValidate}%
  Skriva en återanvändbar funktion som läser in och kontrollerar
  användarens inmatning och frågar om igen tills den duger.
\end{restatable}% Lo08, Lo06
\begin{restatable}{lo}{GuessLOCombine}\label{GuessLOCombine}%
  Utveckla ett program som kombinerar behållare, upprepningar,
  funktioner och felhantering, genom stegvis förfining från överblick
  till körbar kod.
\end{restatable}% Lo02, Lo03
\begin{restatable}{lo}{GuessLOReproduce}\label{GuessLOReproduce}%
  Förklara varför ett program som slumpar är svårt att testa, och göra
  en körning upprepbar genom att sätta slumpgeneratorns startvärde.
\end{restatable}% Lo06

\emph{Förkunskaper}
Modulen förutsätter listor, tupler, \mintinline{python}{for}- och
\mintinline{python}{while}-slingor, egna funktioner samt felhantering
med \mintinline{python}{try}/\mintinline{python}{except}.

\emph{Läsning}
Kursens videogenomgång \emph{Behållare: Gissningsspel} går igenom samma
program.
